\documentclass[10pt,a4paper]{amsart}
\usepackage{amsmath,amssymb,amsthm,graphicx}
\usepackage{verbatim}
\newenvironment{reference}{\comment}{\endcomment}
%\newenvironment{reference}{}{}
\newenvironment{slogan}{\comment}{\endcomment}
\newenvironment{history}{\comment}{\endcomment}

% For commutative diagrams we use Xy-pic
\usepackage[all]{xy}

% We use 2cell for 2-commutative diagrams.
\xyoption{2cell}
% Theorem environments.
%
\theoremstyle{plain}
\newtheorem{theorem}[subsection]{Theorem}
\newtheorem{proposition}[subsection]{Proposition}
\newtheorem{lemma}[subsection]{Lemma}

\theoremstyle{definition}
\newtheorem{definition}[subsection]{Definition}
\newtheorem{example}[subsection]{Example}
\newtheorem{exercise}[subsection]{Exercise}
\newtheorem{situation}[subsection]{Situation}

\theoremstyle{remark}
\newtheorem{remark}[subsection]{Remark}
\newtheorem{remarks}[subsection]{Remarks}

\numberwithin{equation}{subsection}

% Macros
%
\def\lim{\mathop{\mathrm{lim}}\nolimits}
\def\colim{\mathop{\mathrm{colim}}\nolimits}
\def\Spec{\mathop{\mathrm{Spec}}}
\def\Hom{\mathop{\mathrm{Hom}}\nolimits}
\def\Ext{\mathop{\mathrm{Ext}}\nolimits}
\def\SheafHom{\mathop{\mathcal{H}\!\mathit{om}}\nolimits}
\def\SheafExt{\mathop{\mathcal{E}\!\mathit{xt}}\nolimits}
\def\Sch{\mathit{Sch}}
\def\Mor{\mathop{\mathrm{Mor}}\nolimits}
\def\Ob{\mathop{\mathrm{Ob}}\nolimits}
\def\Sh{\mathop{\mathit{Sh}}\nolimits}
\def\NL{\mathop{N\!L}\nolimits}
\def\CH{\mathop{\mathrm{CH}}\nolimits}
\def\proetale{{pro\text{-}\acute{e}tale}}
\def\etale{{\acute{e}tale}}
\def\QCoh{\mathit{QCoh}}
\def\Ker{\mathop{\mathrm{Ker}}}
\def\Im{\mathop{\mathrm{Im}}}
\def\Coker{\mathop{\mathrm{Coker}}}
\def\Coim{\mathop{\mathrm{Coim}}}

% Boxtimes
%
\DeclareMathSymbol{\boxtimes}{\mathbin}{AMSa}{"02}

%
% Macros for moduli stacks/spaces
%
\def\QCohstack{\mathcal{QC}\!\mathit{oh}}
\def\Cohstack{\mathcal{C}\!\mathit{oh}}
\def\Spacesstack{\mathcal{S}\!\mathit{paces}}
\def\Quotfunctor{\mathrm{Quot}}
\def\Hilbfunctor{\mathrm{Hilb}}
\def\Curvesstack{\mathcal{C}\!\mathit{urves}}
\def\Polarizedstack{\mathcal{P}\!\mathit{olarized}}
\def\Complexesstack{\mathcal{C}\!\mathit{omplexes}}
% \Pic is the operator that assigns to X its picard group, usage \Pic(X)
% \Picardstack_{X/B} denotes the Picard stack of X over B
% \Picardfunctor_{X/B} denotes the Picard functor of X over B
\def\Pic{\mathop{\mathrm{Pic}}\nolimits}
\def\Picardstack{\mathcal{P}\!\mathit{ic}}
\def\Picardfunctor{\mathrm{Pic}}
\def\Deformationcategory{\mathcal{D}\!\mathit{ef}}
\usepackage{cleveref}
\crefname{originallemma}{original lemma}{original lemmas}
\Crefname{originallemma}{Original lemma}{Original lemmas}
\makeatletter
\newlabel{algebra-lemma-completion-Noetherian-Noetherian}{{0316}{source}{}{sourcepointer-1}{}}
\newlabel{algebra-lemma-completion-Noetherian-Noetherian@cref}{{[originallemma][0][]0316}{[1][0][]source}{}{}{}}
\makeatother
\crefname{originalreference}{original reference}{original references}
\Crefname{originalreference}{Original reference}{Original references}
\makeatletter
\newlabel{lemma-change-of-fields-smooth}{{07WD}{source}{}{sourcepointer-2}{}}
\newlabel{lemma-change-of-fields-smooth@cref}{{[originalreference][0][]07WD}{[1][0][]source}{}{}{}}
\makeatother
\crefname{originalreference}{original reference}{original references}
\Crefname{originalreference}{Original reference}{Original references}
\makeatletter
\newlabel{lemma-exists-smooth}{{06SM}{source}{}{sourcepointer-3}{}}
\newlabel{lemma-exists-smooth@cref}{{[originalreference][0][]06SM}{[1][0][]source}{}{}{}}
\makeatother
\crefname{originallemma}{original lemma}{original lemmas}
\Crefname{originallemma}{Original lemma}{Original lemmas}
\makeatletter
\newlabel{more-algebra-lemma-formally-smooth}{{07EC}{source}{}{sourcepointer-4}{}}
\newlabel{more-algebra-lemma-formally-smooth@cref}{{[originallemma][0][]07EC}{[1][0][]source}{}{}{}}
\makeatother
\crefname{originallemma}{original lemma}{original lemmas}
\Crefname{originallemma}{Original lemma}{Original lemmas}
\makeatletter
\newlabel{more-algebra-lemma-formally-smooth-completion}{{07ED}{source}{}{sourcepointer-5}{}}
\newlabel{more-algebra-lemma-formally-smooth-completion@cref}{{[originallemma][0][]07ED}{[1][0][]source}{}{}{}}
\makeatother
\crefname{originallemma}{original lemma}{original lemmas}
\Crefname{originallemma}{Original lemma}{Original lemmas}
\makeatletter
\newlabel{more-algebra-lemma-lift-continuous}{{07NJ}{source}{}{sourcepointer-6}{}}
\newlabel{more-algebra-lemma-lift-continuous@cref}{{[originallemma][0][]07NJ}{[1][0][]source}{}{}{}}
\makeatother
\crefname{originalreference}{original reference}{original references}
\Crefname{originalreference}{Original reference}{Original references}
\makeatletter
\newlabel{situation-change-of-fields}{{07W8}{source}{}{sourcepointer-7}{}}
\newlabel{situation-change-of-fields@cref}{{[originalreference][0][]07W8}{[1][0][]source}{}{}{}}
\makeatother
\title{Change of residue field for a versal formal object}
\author{The Stacks Project Authors}
\date{stacks-project:formal-defos.tex:local-snapshot-20261012; GNU FDL 1.2 or later}
\begin{document}
\maketitle
\noindent\textit{Editorial note.} All delivered mathematical text is the author's own, copied from the pinned Stacks Project chapter file formal-defos.tex; no statement or proof step was rewritten, reordered or removed.\par\medskip
\noindent\textit{Editorial note.} The counted proof is the author's complete argument: the construction of S as the adic completion of a polynomial ring over R together with its category and formal smoothness properties, the reduction to a smoothness statement via the change of residue field lemma, and the verification of the lifting property on morphism sets.\par\medskip
\noindent\textit{Editorial note.} The label lemma-change-of-field-versal-ring names the counted statement, and the proof environment immediately following it in the source is delivered verbatim as its proof.\par\medskip
\noindent\textit{Editorial note.} The counted statement is carried by the author's own lemma environment in the chapter source.\par\medskip
\noindent\textit{Editorial note.} Pre-open gate: python operations/agent-authoring/tools/precard\_gate.py --source formal-defos.tex --label lemma-change-of-field-versal-ring --start 7692 --end 7753 --local-id 28017 returned RESULT=OK (L28017 free, occupancy FREE, public interval clear). Candidate taken from sources/stacks-formal-defos/candidates-halfA.json.\par\medskip
\noindent\textit{Editorial note.} The statement and proof are copied from the pinned chapter file; the cited inputs are kept as external pointers with their canonical Stacks tags.\par\medskip
\noindent\textit{Allowed standard prerequisites (editorial).}\begin{itemize}
\item Smoothness of the transformation induced by a change of residue field. Same chapter, the lemma that a smooth transformation stays smooth after the change of residue field construction. Supplies the smoothness of the induced transformation from the localized functor to the changed functor.
\item Formal smoothness of a power series over a Noetherian ring in the adic topology. More on Algebra, the lemmas identifying formal smoothness of a ring map with the lifting property and its compatibility with completion. Supplies the formal smoothness of R -> S in the m\_S-adic topology and the existence of the lift in the last step.
\item Membership of the completion in the category of Artinian quotients. Algebra, the lemma that the completion of a Noetherian local ring is again Noetherian and has the expected residue field. Supplies the membership of S in the category over l.
\end{itemize}\medskip
\ifdefined\hypertarget\hypertarget{sourcepointer-1}{}\fi
\noindent\textit{Source reference.} Original lemma 0316: sources/stacks-chow/tags.txt, source lines 3973--3973. Prerequisite pointer invoked at proof line(s) 7711; it names a standard fact of a sibling chapter that the counted argument uses, and every step of the counted proof built on it is retained verbatim in the delivered body.\par\medskip
\ifdefined\hypertarget\hypertarget{sourcepointer-2}{}\fi
\noindent\textit{Source reference.} Original reference 07WD: sources/stacks-formal-defos/pinned-original-source.tex, source lines 7674--7674. Prerequisite pointer invoked at proof line(s) 7722; it names a standard fact of the same chapter that the counted argument uses, and every step of the counted proof built on it is retained verbatim in the delivered body.\par\medskip
\ifdefined\hypertarget\hypertarget{sourcepointer-3}{}\fi
\noindent\textit{Source reference.} Original reference 06SM: sources/stacks-formal-defos/pinned-original-source.tex, source lines 2598--2598. Prerequisite pointer invoked at proof line(s) 7717; it names a standard fact of the same chapter that the counted argument uses, and every step of the counted proof built on it is retained verbatim in the delivered body.\par\medskip
\ifdefined\hypertarget\hypertarget{sourcepointer-4}{}\fi
\noindent\textit{Source reference.} Original lemma 07EC: sources/stacks-chow/tags.txt, source lines 9442--9442. Prerequisite pointer invoked at proof line(s) 7714; it names a standard fact of a sibling chapter that the counted argument uses, and every step of the counted proof built on it is retained verbatim in the delivered body.\par\medskip
\ifdefined\hypertarget\hypertarget{sourcepointer-5}{}\fi
\noindent\textit{Source reference.} Original lemma 07ED: sources/stacks-chow/tags.txt, source lines 9443--9443. Prerequisite pointer invoked at proof line(s) 7715; it names a standard fact of a sibling chapter that the counted argument uses, and every step of the counted proof built on it is retained verbatim in the delivered body.\par\medskip
\ifdefined\hypertarget\hypertarget{sourcepointer-6}{}\fi
\noindent\textit{Source reference.} Original lemma 07NJ: sources/stacks-chow/tags.txt, source lines 9764--9764. Prerequisite pointer invoked at proof line(s) 7752; it names a standard fact of a sibling chapter that the counted argument uses, and every step of the counted proof built on it is retained verbatim in the delivered body.\par\medskip
\ifdefined\hypertarget\hypertarget{sourcepointer-7}{}\fi
\noindent\textit{Source reference.} Original reference 07W8: sources/stacks-formal-defos/pinned-original-source.tex, source lines 7494--7494. Prerequisite pointer invoked at proof line(s) 7694; it names a standard fact of the same chapter that the counted argument uses, and every step of the counted proof built on it is retained verbatim in the delivered body.\par\medskip
\begin{lemma}
\label{lemma-change-of-field-versal-ring}
With notation and assumptions as in Situation \ref{situation-change-of-fields}.
Let $\xi$ be a versal formal object for $\mathcal{F}$ lying over
$R \in \Ob(\widehat{\mathcal{C}}_{\Lambda, k})$. Then there exist
\begin{enumerate}
\item an $S \in \Ob(\widehat{\mathcal{C}}_{\Lambda, l})$
and a local $\Lambda$-algebra homomorphism $R \to S$ which is
formally smooth in the $\mathfrak m_S$-adic topology and induces
the given field extension $l/k$ on residue fields, and
\item a versal formal object of $\mathcal{F}_{l/k}$ lying over $S$.
\end{enumerate}
\end{lemma}

\begin{proof}
Construction of $S$. Choose a surjection $R[x_1, \ldots, x_n] \to l$
of $R$-algebras. The kernel is a maximal ideal $\mathfrak m$.
Set $S$ equal to the $\mathfrak m$-adic completion of the Noetherian ring
$R[x_1, \ldots, x_n]$. Then $S$ is in $\widehat{\mathcal{C}}_{\Lambda, l}$
by Algebra, Lemma \ref{algebra-lemma-completion-Noetherian-Noetherian}.
The map $R \to S$ is formally smooth in the $\mathfrak m_S$-adic topology by
More on Algebra, Lemmas
\ref{more-algebra-lemma-formally-smooth} and
\ref{more-algebra-lemma-formally-smooth-completion}
and the fact that $R \to R[x_1, \ldots, x_n]$ is formally smooth.
(Compare with the proof Lemma \ref{lemma-exists-smooth}.)

\medskip\noindent
Since $\xi$ is versal, the transformation
$\underline{\xi} : \underline{R}|_{\mathcal{C}_{\Lambda, k}} \to \mathcal{F}$
is smooth. By Lemma \ref{lemma-change-of-fields-smooth} the induced map
$$
(\underline{R}|_{\mathcal{C}_{\Lambda, k}})_{l/k}
\longrightarrow
\mathcal{F}_{l/k}
$$
is smooth. Thus it suffices to construct a smooth morphism
$\underline{S}|_{\mathcal{C}_{\Lambda, l}} \to
(\underline{R}|_{\mathcal{C}_{\Lambda, k}})_{l/k}$.
To give such a map means for every object $B$ of $\mathcal{C}_{\Lambda, l}$
a map of sets
$$
\Mor_{\widehat{\mathcal{C}}_{\Lambda, l}}(S, B)
\longrightarrow
\Mor_{\widehat{\mathcal{C}}_{\Lambda, k}}(R, B \times_l k)
$$
functorial in $B$. Given an element $\varphi : S \to B$
on the left hand side we send it to the composition $R \to S \to B$
whose image is contained in the sub $\Lambda$-algebra $B \times_l k$.
Smoothness of the map means that given a surjection
$B' \to B$ and a commutative diagram
$$
\xymatrix{
S \ar[r] & B \ar@{=}[r] & B \\
R \ar[u] \ar[r] & B' \times_l k \ar[u] \ar[r] & B' \ar[u]
}
$$
we have to find a ring map $S \to B'$ fitting into the outer rectangle.
The existence of this map is guaranteed as we chose $R \to S$ to be
formally smooth in the $\mathfrak m_S$-adic topology, see
More on Algebra, Lemma \ref{more-algebra-lemma-lift-continuous}.
\end{proof}

\end{document}
